% !TeX program = pdflatex
\documentclass[11pt]{article}
\usepackage{amsmath,amssymb,amsthm}
\usepackage[margin=28mm]{geometry}
\newtheorem{theorem}{Theorem}
\title{Spectral stability from the variational principle}
\author{}\date{}
\begin{document}
\maketitle
\section{Operator norm}
For $A\in\mathbb C^{m\times n}$, write
$\|A\|_{\mathrm{op}}=\sup_{\|z\|_2=1}\|Az\|_2$.
The triangle inequality gives
$\|A\|_{\mathrm{op}}\leq\|B\|_{\mathrm{op}}+\|A-B\|_{\mathrm{op}}$.
Interchanging $A$ and $B$ proves the Lipschitz bound.
\section{A rank-one perturbation}
Let $u$ and $v$ be unit vectors.
Take $\theta>0$ and set $X=\theta uv^*+E$.
The nonzero singular value of $\theta uv^*$ is $\theta$.
\begin{equation}\label{eq:bound}
  |s_1(X)-\theta|\leq\|E\|_{\mathrm{op}}.
\end{equation}
Equation~\eqref{eq:bound} is deterministic.
\section{Variational proof}
For every unit vector $z$,
\[\|Xz\|_2\leq\theta|v^*z|+\|Ez\|_2
  \leq\theta+\|E\|_{\mathrm{op}}.\]
Taking the supremum proves the upper bound.
For the lower bound, evaluate the bilinear form at $u,v$:
\[s_1(X)\geq|u^*Xv|=|\theta+u^*Ev|
  \geq\theta-\|E\|_{\mathrm{op}}.\]
These inequalities hold for every rectangular noise matrix.
\section{Interpretation}
Noise sets the estimation error.
The bound uses no independence or distributional assumption.
A probabilistic model is needed only to control the noise norm.
\end{document}
